% ============================================
% Johns Hopkins Letter Template
% Assistant Professor Cover Letter – Texas State University (McCoy College of Business)
% ============================================

\documentclass[11pt,letterpaper]{letter}

\usepackage{graphicx}
\usepackage{eso-pic}
\usepackage{geometry}
\usepackage{microtype}
\usepackage[hidelinks]{hyperref}

% ----- Page Setup -----
\geometry{left=25mm,right=25mm,top=25mm,bottom=25mm}

% ----- Logo Placement (foreground, first page only) -----
\newcommand{\PutJHULogoFG}[1]{%
  \AddToShipoutPictureFG*{%
    \AtPageUpperLeft{%
      \hspace{10mm}%
      \raisebox{-38mm}{%
        \includegraphics[width=6cm]{#1}%
      }%
    }%
  }%
}

% ----- Sender Information -----
\signature{Decory Edwards}
\address{Department of Economics \\ Johns Hopkins University \\ Baltimore, MD 21218 \\ \texttt{dedwar65@jh.edu}}

\begin{document}
\PutJHULogoFG{Images/jhulogo.png}

\begin{letter}{Search Committee \\
Department of Finance \& Economics \\
McCoy College of Business, Texas State University \\
McCoy Hall 504, 601 University Drive \\
San Marcos, TX 78666-4616, USA}

\opening{Dear Members of the Search Committee,}

I am writing to express my interest in the Assistant Professor position in Economics (macroeconomics) in the Department of Finance \& Economics at Texas State University’s McCoy College of Business. I am a Ph.D.\ candidate in Economics at Johns Hopkins University, advised by Professors Christopher D.\ Carroll, Nicholas W.\ Papageorge, and Stelios Fourakis, and I expect to receive my degree in May 2026. My research fields are macroeconomics, wealth and income dynamics, and computational economics.

My research agenda is unified by a focus on understanding the determinants of wealth inequality and the mechanisms through which heterogeneity in returns to wealth shapes aggregate outcomes. In my job market paper, I estimate the degree of heterogeneity in individual portfolio returns using micro data from the Survey of Consumer Finances and simulate a structural model in which households face idiosyncratic return risk. I show that allowing for persistent heterogeneity in returns substantially improves the model’s ability to match the empirical distribution of wealth, offering a tractable way to connect micro-level return dispersion to macroeconomic inequality.

In a second project, I use the Health and Retirement Study to study the relationship between interpersonal trust and portfolio performance. Building on the literature that finds a hump-shaped relationship between trust and income, I show that a similar relationship holds for individual returns to wealth measured in the HRS data, highlighting a behavioral channel through which social attitudes shape saving and investment outcomes. This work also provides new estimates of the persistent component of returns to net worth, which motivate the calibration of heterogeneity in my structural model.

\textbf{Fit with McCoy and the position.} The business-school setting at McCoy is an excellent match for my empirical and computational approach. My work links household portfolio behavior and return heterogeneity to questions central to business and policy analysis—capital allocation under risk, the distributional impact of taxation, and the interpretation of micro data for forecasting and decision support. I welcome the opportunity to contribute across McCoy’s programs and centers while pursuing rigorous, policy-relevant research.

\textbf{Teaching.} My teaching experience centers on undergraduate macroeconomics and an upper-level macro policy seminar, where I have led weekly discussion sections and held regular office hours for classes ranging from 16 to 40 students. Consistent with my teaching statement, my approach is student-centered: I aim to help students ``convince themselves'' that they have moved from confusion to understanding by pairing clear scaffolding with structured office-hour coaching that builds trust and confidence. At Texas State, I am prepared to teach macroeconomics at the undergraduate and graduate levels and contribute to related quantitative and analytics courses; I am fully comfortable teaching at both the San Marcos and Round Rock campuses. I would also welcome developing electives such as \textit{Economics of Inequality and Tax Policy} or \textit{Computational Economics for Business Decisions}, emphasizing reproducible workflows and evidence-based decision-making.

I believe I would be a strong fit because my computational and empirical toolkit allows me to (i) produce robust, data-backed estimates of returns and distributional outcomes, (ii) build and validate models that speak to business, regulatory, and policy-relevant questions at the intersection of macroeconomics and strategy, and (iii) collaborate across teams to present insights in concise, actionable formats for diverse audiences. I am excited by McCoy’s applied, collaborative environment, where academic rigor meets high-impact teaching, research, and community engagement.

I would be honored to contribute to Texas State’s research and teaching mission and to mentor students pursuing quantitative careers in economics and business. Thank you for your consideration; I welcome the opportunity to discuss my fit with McCoy College of Business at your convenience.

\closing{Sincerely,}

\end{letter}
\end{document}
